\usepackage[most]{tcolorbox}
%\usepackage{pgf,tikz}
\usepackage{tkz-base,tkz-fct,tkz-euclide,tkz-tab}
\usetikzlibrary{babel}
\usepackage{pas-math}
\usepackage{tablists}

 % lining
%\newfontfamily{\LinuxLibertine}[Ligatures=TeX]{Linux Libertine O} % ,Numbers=OldStyle
%\newfontfamily{\LinuxBiolinumO}[Ligatures=TeX]{Linux Biolinum O}% ,Numbers=OldStyle
%\newfontfamily{\Ubuntu}[Ligatures=TeX]{Ubuntu} % ,Numbers=OldStyle
%\setsansfont[Color=black]{Arial}
%\setmonofont[Color=black]{Courier New}
%\usepackage{tabularray}
\usepackage{pifont,tasks}
\DeclareMathOperator{\bary}{bar}
 \newcommand{\activity}[3]{%
 	\vspace*{-0.3cm}
     \begin{bclogo}[logo =\bcbook, noborder = true , sousTitre={\bf #1},marge =0,barre=snake,tailleOndu=2.5,epBarre=2]{\fbox{Activit\'e #2} : } %
         #3 \end{bclogo}
     \vspace*{-0.5cm}
    }
     \newcommand{\consigne}[3]{%
     	%\vspace*{-0.2cm}
         \begin{bclogo}[logo =\bccrayon, noborder = true , sousTitre={\small\bf #1},marge =0,epBarre=2,couleurBarre=blue]{Consigne #2 : } %
             #3 \end{bclogo}
         \vspace*{-0.3cm}
        }
         \newcommand{\info}[1]{
         %	\vspace*{-0.5cm}
             \begin{bclogo}[logo =\bcinfo, noborder = true , sousTitre=\textbf{\textsf{\uuline{Information}}} : ,marge =0,motifBarre=$\star$,epBarre=2]\textit{#1} %
                \end{bclogo}
            %\vspace*{-0.5cm}
            }
            
                         % \newcommand{\remark}[1]{%
                         %     \begin{bclogo}[logo =\bclampe, noborder = true,marge =0,epBarre=2]{Remarque : } %
                          %        \textit{#1} \end{bclogo}
                          %  }
                            \newcommand{\resultat}[1]{%
                            	\vspace*{-0.3cm}
                                \begin{bclogo}[logo =\bcfleur, noborder = true,marge =0,epBarre=2,couleurBarre=red]{\uuline{R\'esultats attendus} : } %
                                    #1 \end{bclogo}
                           \vspace*{-0.5cm} }             
                            \newcommand{\exercic}[1]{%
                                \begin{bclogo}[logo =\bcpanchant, noborder = true,marge =0,epBarre=2,couleurBarre=red]{\uuline{Exercices de maison} : } %
                                    #1 \end{bclogo}
                            } 
                            
                             \DefineNewBoxLikeRem{regles}{R\`egles de repr\'esentation en perspective cavali\`ere}{green}{white}{phv}
                             
                              \DefineNewBoxLikeDef{retenons}{Retenons}{green}{white}{phv}
                           
                            
                            \newcommand\V[1]{\overrightarrow{#1}}
                            \newcommand{\F}[2]{\dfrac{#1}{#2}}
                            \newcommand{\f}[2]{\frac{#1}{#2}}
                            \newcommand{\be}{\begin{enumerate}}
                                \newcommand{\ee}{\end{enumerate}}
%%%%%%%%%%%%%%%%%%%%Propriété avec titre%%%%%%%%%%%%%%%%%%%%%%%%%%%%%%

%%%%%%%%%%%%%%%%%%%%%%%Remarque%%%%%%%%%%%%%%%%%%%%
\newtcolorbox{remark}[1][]{enhanced,
	before skip=2mm,after skip=3mm,left skip=6.5mm,
	boxrule=0.4pt,left=4mm,right=2mm,top=1mm,bottom=1mm,
	colback=yellow!50,colframe=yellow!20!black,
	sharp corners,underlay={%
		%	\path[fill=yellow!80!black] ([yshift=3mm]interior.south east)--++(-0.4,-0.1)--++(0.1,-0.2);
		%	\path[draw=blue,shorten <=-0.05mm,shorten >=-0.05mm] ([yshift=3mm]interior.south east)--++(-0.4,-0.1)--++(0.1,-0.2);
		\path[fill=yellow!50!black,draw=none] (interior.south west) rectangle node[white]{\huge\bfseries !} ([xshift=4mm]interior.north west);
	},width=\linewidth,drop fuzzy shadow,#1}
%%%%%%%%%%%%%%%%%%%%%%%%%Propriétés au pluriel%%%%%%%%%%%%%%%%%%
\newtcolorbox{proprietys}[1][]{enhanced,
	before skip=2mm,after skip=3mm,
	boxrule=0.4pt,left=10mm,right=2mm,top=1mm,bottom=1mm,
	colback=yellow!50,colframe=yellow!20!black,
	sharp corners,underlay={%
		%	\path[fill=yellow!80!black] ([yshift=3mm]interior.south east)--++(-0.4,-0.1)--++(0.1,-0.2);
		%	\path[draw=blue,shorten <=-0.05mm,shorten >=-0.05mm] ([yshift=3mm]interior.south east)--++(-0.4,-0.1)--++(0.1,-0.2);
		\path[fill=yellow!50!black,draw=none] (interior.south west) rectangle node[white,rotate=90]{\huge\bfseries Propriétés} ([xshift=10mm]interior.north west);
	},drop fuzzy shadow,#1}
%%%%%%%%%%%%%%%%%%%%%%%%%%%%%%%%%%%%%%%%%%%%%%%%%%%%%%%%%%%%%%%%%%%%%


%%%%%%%%%%%%%%%%%%%%%%%%SITUATION DE DÉPART %%%%%%%%%%%%%%%%%%%%%%%
          		\newtcbox{\mylib}{enhanced,nobeforeafter,tcbox raise base,boxrule=0.4pt,top=0mm,bottom=0mm,
	right=0mm,left=4mm,arc=1pt,boxsep=3pt,before upper={\vphantom{dlg}},
	colframe=green!50!black,coltext=green!25!black,colback=white,
	overlay={\begin{tcbclipinterior}\fill[green!75!blue!50!white] (frame.south west)
			rectangle node[text=white,font=\sffamily\bfseries\tiny,rotate=90] {S.A 1} ([xshift=4mm]frame.north west);\end{tcbclipinterior}}}
		
%%%%%%%%%%%%%%%%%%%%%%%%%%%%%%%%%%%%%%%%%%%%ù%%%%%%%%%%%%%%%%%%%%%%ù%%%%%%%%%%%%%%%%%%%%%%%%%%%%%%%%%%%%%%%%%%%%%%%%%%%%%%%%%%%%ù%%%%%%%%%%%%%%%%ù  Commandes mathématiques%%%%%%%%%%%%%%%%
%%%%%%%%%%%%%%%%%%%%%%%%%%%%%%%%%%%%%%%%%%%%%%%%%%%%%%%%%%%%%%%


\DeclareRobustCommand{\ve}[1]{%
	\overrightarrow{\rule{0em}{1.8ex} #1 \rule{0.15em}{0ex}}}
% %%% Commande pour dï¿œaler les parenthï¿œes vers le bas autour
% %%% des vecteurs par exemple :
\newcommand{\DecalV}[2][2pt]{%
	\raisebox{#1}{%
		$\left(\raisebox{-#1}{\ensuremath{#2}}\right)$}}
%%%%angle de vecteur avec chapeau
\newcommand{\anglevec}[2]{%
	\DecalV{%
		\widehat{%
			\ve{#1};\ve{#2}
}}}

%\newcommand{\I}{{\rm i}}

%\newtcolorbox{encadre}{fonttitle=\scriptsize,colback=LightBlue!10!white,colframe=blue,width=4.7cm,valign=center,halign=left}

%%%%%%%%%%%%%%%%%%%%%%%%%%%%%%%%%%
% Definition de la date : les mois sont écrits en français
% Date Mois Année
\newcommand{\Date}{
	%\ifnum\number\day<10 0\fi \number\day \space%
	\ifcase \month \or Janvier \or Février \or Mars \or Avril  \or Mai%
	\or Juin\or Juillet\or Août \or Septembre \or Octobre \or Novembre \or Décembre \fi%
	\space \number \year}

% Commande \longarc pour afficher les arc orientés
\newcommand{\ArcOriente}[1]{\ensuremath{%
		\begin{tikzpicture}[baseline,anchor=base,inner sep=0pt,outer sep=0pt,line width=0.5pt]%
		\node at (0,0) {$#1$};%
		\draw[<-,>=stealth',yshift=-1pt,xshift=-0.04cm] (45:0.5cm) arc (45:135:0.4cm);%
		\end{tikzpicture}%
	}%
}

% Commande \longarc pour afficher les arc orientés
\newcommand{\Arc}[1]{\ensuremath{%
		\begin{tikzpicture}[baseline,anchor=base,inner sep=0pt,outer sep=0pt,line width=0.5pt]%
		\node at (0,0) {$#1$};%
		\draw[yshift=-1pt,xshift=-0.04cm] (45:0.5cm) arc (45:135:0.4cm);%
		\end{tikzpicture}%
	}%
}
%%%%%%%%%%%%%%%%%%%%%%%%%%%%%%%
%   Symbole parallel incliné                                                                      %
%%%%%%%%%%%%%%%%%%%%%%%%%%%%%%%%%%%%%%%%%%%%%%%%%%%%%%%%%%%%%%%%%%%%%%%%%%%%%%%%%%%%%%%%%%%%%%%%%%%
\newcommand{\parallels}{\Scale[1.05]{\sslash}}
\newcommand{\parallelsl}{\mathbin{\!/\mkern-5mu/\!}}


\newcommand{\segment}[1]{\ensuremath{\big[#1\big]}}
\newcommand{\Segment}[1]{\ensuremath{\left[#1\right]}}
\newcommand{\droite}[1]{\ensuremath{\big(#1\big)}}
\newcommand{\Droite}[1]{\ensuremath{\left(#1\right)}}
\newcommand{\plan}[1]{\ensuremath{\big(#1\big)}}
\newcommand{\Plan}[1]{\ensuremath{\left(#1\right)}}
\newcommand{\ortho}{\Scale[1.2]{\bot}}
\newcommand{\scalaire}[2]{\ensuremath{\vv{#1}{\hbox{\fontsize{18}{18}\selectfont$\mathbf{.}$}\vv{#2}}}}
\newcommand{\module}[1]{\ensuremath{\big|#1\big|}}
% commande \reallywidehat : mesure angle orienté
% necessite: listofitems.sty , listofitems.tex , scalerel.sty et stackengine.sty
%\stackMath
%\newcommand\reallywidehat[1]{%
%	\savestack{\tmpbox}{\stretchto{%
%			\scaleto{%
%				\scalerel*[\widthof{\ensuremath{#1}}]{\kern4pt\bigwedge\kern4pt}%
				%{\rule[-\textheight/2]{1ex}{\textheight}}%WIDTH-LIMITED BIG WEDGE
%			}{\textheight}%
%		}{0.5ex}}%
%	\stackon[2pt]{#1}{\tmpbox}%
%}

% commande \mao{#1}{#2} : mesure angle orienté
\newcommand{\mao}[2]{\ensuremath{
		{\fontsize{3}{3}\selectfont\ }\reallywidehat{\left(\vv{#1},\vv{#2}\right)}}
} 
%
%%%%%%%%%%%%%%%%%%%%%%%%%%%%%%%%%%%%%%%%%%%%%%%%%%%%%%%%%%%%%%%%%%%%%%%%%%%%%%%%%%%%%%%%%%%%%%%%%%%		
		
\definecolor{berry}{RGB}{235, 37, 129}
%%%%%%%%%%%%%%%%%%%%%%%%%%%
\newmdenv[skipabove=7pt,
skipbelow=7pt,
rightline=false,
leftline=false,
topline=false,
bottomline=false,
linecolor=berry,
backgroundcolor=berry!10,
innerleftmargin=5pt,
innerrightmargin=5pt,
innertopmargin=2pt,
leftmargin=0cm,
rightmargin=0cm,
linewidth=2.5pt,
innerbottommargin=5pt
]{boite}
%%%%%%%%%%%%%%%%%%%%%%%%%%%%